Throughout this section let \(Z_1,\dots,Z_n\) be i.i.d.\ \(\cN(0,I_2)\), let
\(\rho_n=\sqrt{2\log n}\), set \(\bar Z_n=n^{-1}\sum_iZ_i\), and put
\[
    H_n(t)=\frac1n\sum_{i=1}^n\exp\{t\cdot Z_i-|t|^2/2\},
    \qquad
    D_n(t)=e^{-t\cdot\bar Z_n}e^{-\kappa_n|t|^2}H_n(t),
    \qquad
    \lambda_n=\kappa_n\rho_n^2 .
\]
The one-atom event is \(E_n^{(1)}=\{\sup_tD_n(t)\le1\}\).
We write \(\Phi\) and \(\phi\) for the standard normal distribution
function and density.

\subsection{The limiting profile}

The KKT supremum chooses its direction after seeing the sample. A fixed
angular interval of width \(1/\rho_n\) contains a point with bounded radial
excess with probability \(O(1/\rho_n)\), but the maximizing interval can be
centered on such a point after the sample is observed. This selected extreme
creates one random spike. The many observations whose radial energies lie
just below \(\log n\) produce the deterministic background.

Write
\[
    Z_i=R_iU_i,\qquad U_i=e^{i\Theta_i}\in S^1,\qquad
    \xi_i=\frac{R_i^2}{2}-\log n ,
\]
and let \(\operatorname{ang}(\phi-\psi)\in(-\pi,\pi]\) denote the principal
signed angular displacement from direction \(\psi\) to direction \(\phi\). We
also write \(\operatorname{ang}(u,v)\) for the same displacement from \(v\) to
\(u\), and set \(Y_i(v)=\rho_n\operatorname{ang}(U_i,v)\). The global radial edge process
\[
    \Xi_n=\sum_{i=1}^n\delta_{(\xi_i,\Theta_i)}
\]
converges to a Poisson point process \(\Xi=\sum_j\delta_{(\xi_j,\Theta_j)}\)
on \(\R\times\mathbb T\), where \(\mathbb T=\R/2\pi\mathbb Z\), with intensity
\[
    e^{-x}\,dx\,\frac{d\phi}{2\pi}.
\]
For a fixed deterministic direction \(\psi\), however, the locally rescaled
process
\[
    \sum_i
    \delta_{\left(\xi_i,\ \rho_n\operatorname{ang}(\Theta_i-\psi)\right)}
\]
has intensity
\[
    \frac{1+o(1)}{2\pi\rho_n}\,e^{-x}\,dx\,dy
\]
on compact \(x,y\)-sets, and hence converges to the empty process. The
nontrivial local limit is obtained only after conditioning on, and recentering
around, a radial edge point.

For \(t=(\rho_n+s)e^{i\psi}\), one observation contributes
\[
    W_{n,i}(\psi,s)
    =
    \frac1n
    \exp\left\{(\rho_n+s)e^{i\psi}\cdot Z_i
        -\frac{(\rho_n+s)^2}{2}\right\}.
\]
If \(\xi_i\), \(s\), and
\[
    y_i=\rho_n\operatorname{ang}(\Theta_i-\psi)
\]
are bounded, then
\[
    W_{n,i}(\psi,s)
    =
    \exp\{\xi_i-s^2/2-y_i^2/2+o(1)\}.
\]
Thus a radial edge point \((\xi_j,\Theta_j)\), viewed in the angular coordinate
\(\psi=\Theta_j+y/\rho_n\), creates the local spike
\[
    e^{\xi_j-s^2/2-y^2/2}.
\]
The spike has \(O(1)\) width in blown-up coordinates, hence
\(O(1/\rho_n)\) width on the original circle.

For fixed \(A,C<\infty\), the \(O(1)\)-radial edge spikes are separated:
\[
    \Pbb\left\{
    \exists i\ne j:\ \xi_i,\xi_j\ge-A,\
    \rho_n|\operatorname{ang}(\Theta_i-\Theta_j)|\le C
    \right\}\to0 ,
\]
because the expected number of such pairs is
\[
    O\left\{
        \left(\int_{-A}^{\infty}e^{-x}\,dx\right)^2
        \frac C{\rho_n}
    \right\}=o(1).
\]
Consequently, after the uniform compactness estimates are proved, at most one
radial extreme contributes at order one near any maximizing
direction. All other observations form the deterministic background. The
limiting event is therefore determined by the largest radial excess rather
than by a nontrivial random process on the entire circle.

The following calculation identifies the background. Let
\(t=(\rho_n+s)v\), \(v\in S^1\), and truncate observations at the
one-dimensional projection edge in direction \(v\):
\[
    H_n^{\mathrm{diff}}(v,s)
    =
    \frac1n\sum_{i=1}^n
    \exp\left\{(\rho_n+s)v\cdot Z_i-\frac{(\rho_n+s)^2}{2}\right\}
    \mathbf 1_{\{v\cdot Z_i\le\rho_n\}}
\]
satisfies
\[
    \E H_n^{\mathrm{diff}}(v,s)=\Phi(-s),
    \qquad
    \operatorname{Var}H_n^{\mathrm{diff}}(v,s)
    \le
    \frac1n e^{(\rho_n+s)^2}\Phi(-\rho_n-2s)
    \sim
    \frac{e^{-s^2}}{\sqrt{2\pi}\rho_n}
\]
uniformly for \(s\) in compact sets.  Thus the non-extreme cloud contributes \(\Phi(-s)\) to \(H_n\), or \(e^{-\lambda}\Phi(-s)\) after the damping factor in \(D_n\).

A radial edge point with
\[
    \xi_i=\frac{|Z_i|^2}{2}-\log n
\]
and direction \(U_i\) contributes, at \(t=(\rho_n+s)U_i\),
\[
    e^{-\lambda}\exp\{\xi_i-s^2/2+o(1)\}.
\]
Combining the two contributions, a radial excess \(x\) can produce a crossing
precisely when
\[
    g(x):=\sup_{s\in\R}\{\Phi(-s)+e^{x-s^2/2}\}
\]
exceeds the damped threshold \(e^\lambda\). This gives the one-atom limit
\[
    Q_2(\lambda)=\exp\{-e^{-x_\lambda}\},
\]
where \(x_\lambda\) is uniquely determined by
\[
    \sup_{s\in\R}\{\Phi(-s)+e^{x_\lambda-s^2/2}\}=e^\lambda .
\]
Equivalently, if \(a_\lambda>0\) solves
\[
    \Phi(a_\lambda)+\frac{\phi(a_\lambda)}{a_\lambda}=e^\lambda,
\]
then
\[
    x_\lambda=-\log(\sqrt{2\pi}a_\lambda),
    \qquad
    Q_2(\lambda)=\exp\{-\sqrt{2\pi}a_\lambda\}.
\]
The optimizer is \(s=-a_\lambda\), obtained from
\[
    -\phi(s)-s e^{x-s^2/2}=0.
\]
The map \(a\mapsto \Phi(a)+\phi(a)/a\) is strictly decreasing, since its derivative is \(-\phi(a)/a^2\), so \(a_\lambda\) is unique.

The optimized profile crosses every level above one exactly once. Its strict
monotonicity later converts a fixed radial gap into a fixed KKT gap.
\begin{lemma}
\label{lem:q-two-curve}
Let
\[
    g(x)=\sup_{s\in\R}\{\Phi(-s)+e^{x-s^2/2}\}.
\]
Then \(g\) is finite, continuous, and strictly increasing, with
\(\lim_{x\to-\infty}g(x)=1\) and \(\lim_{x\to\infty}g(x)=\infty\).  Hence,
for every \(\lambda>0\), there is a unique \(x_\lambda\) such that
\(g(x_\lambda)=e^\lambda\), and for every \(\eta>0\),
\[
    g(x_\lambda-\eta)<e^\lambda<g(x_\lambda+\eta).
\]
\end{lemma}

\begin{proof}
For every finite \(x\), the objective exceeds \(1\) at some finite negative
\(s\). Indeed, with \(s=-a\), Mills' ratio gives
\(1-\Phi(a)\sim\phi(a)/a\), whereas
\(e^{x-a^2/2}=\sqrt{2\pi}e^x\phi(a)\); the second term dominates the first
for sufficiently large \(a\). On a compact \(x\)-interval, the objective
converges uniformly to \(1\) as \(s\to-\infty\) and to \(0\) as
\(s\to+\infty\). Its supremum is therefore attained in a common compact
\(s\)-interval. Continuity follows from compact maximization, and strict
monotonicity follows because \(e^{x-s^2/2}\) is strictly increasing in \(x\)
at every finite maximizer. Finally,
\(1\le g(x)\le1+e^x\) implies \(g(x)\to1\) as \(x\to-\infty\), while
\(g(x)\ge e^x\to\infty\) as \(x\to\infty\).
\end{proof}

The explicit parametrization also gives the endpoint behavior
\[
    1-Q_2(\lambda)\sim e^{-\lambda}\quad(\lambda\to\infty),
    \qquad
    \log Q_2(\lambda)
    =
    -2\sqrt{\pi\log(1/\lambda)}\{1+o(1)\}
    \quad(\lambda\downarrow0).
\]
Indeed, as \(\lambda\to\infty\), the defining equation gives
\(a_\lambda\sim\phi(0)e^{-\lambda}\). Since
\(e^{-x_\lambda}=\sqrt{2\pi}a_\lambda\), this proves the first relation.
As \(\lambda\downarrow0\), the two-term Mills expansion gives
\[
 e^\lambda-1
 =\frac{\phi(a_\lambda)}{a_\lambda^3}
   \{1+O(a_\lambda^{-2})\}.
\]
Thus \(a_\lambda^2\sim2\log(1/\lambda)\), which proves the second relation.
For large \(\lambda\), the diffuse background is negligible and the radial
Gumbel tail determines the limit. By contrast,
\(Q_2(\lambda)\downarrow0\) as the shrinkage tends to zero.

\subsection{Radial compactness}

Radial compactness reduces all possible crossings to
\(|t|=\rho_n+O(1)\), without yet identifying the limiting field.

\begin{proposition}
\label{prop:d-two-radial-compactness}
Assume \(d=2\) and \(\lambda_n=\kappa_n\rho_n^2\to\lambda\in(0,\infty)\),
where \(\rho_n=\sqrt{2\log n}\).  Then
\[
    \lim_{A\to\infty}\limsup_{n\to\infty}
    \Pbb\left\{
        \sup_{\left||t|-\rho_n\right|>A}D_n(t)>1
    \right\}=0 .
\]
Thus, up to a probability which vanishes after \(n\to\infty\) and then
\(A\to\infty\), every crossing has the form
\[
    t=(\rho_n+s)v,\qquad |s|\le A,\quad v\in S^1 .
\]
\end{proposition}

Write
\[
    H_n^{\mathrm{diff}}(r,v)=\frac1n\sum_i e^{rv\cdot Z_i-r^2/2}
        \mathbf 1_{\{v\cdot Z_i\le \rho_n\}},
    \qquad
    T_n(r,v)=\frac1n\sum_i e^{rv\cdot Z_i-r^2/2}
        \mathbf 1_{\{v\cdot Z_i> \rho_n\}} .
\]
Thus \(H_n(rv)=H_n^{\mathrm{diff}}(r,v)+T_n(r,v)\).  The split is at the
one-dimensional edge in direction \(v\).

The diffuse part cannot create a crossing below the compact window or at its
upper boundary.
\begin{lemma}
\label{lem:d-two-diffuse-outside}
For every \(\delta>0\),
\[
\begin{aligned}
    &\lim_{A\to\infty}\limsup_{n\to\infty}
    \Pbb\left\{
        \sup_{\rho_n/2\le r\le \rho_n-A,\ v\in S^1}
        H_n^{\mathrm{diff}}(r,v)>1+\delta
    \right\}=0,\\
    &\lim_{A\to\infty}\limsup_{n\to\infty}
    \Pbb\left\{
        \sup_{v\in S^1}H_n^{\mathrm{diff}}(\rho_n+A,v)>\delta
    \right\}=0 .
\end{aligned}
\]
\end{lemma}

\begin{proof}
Put
\[
    \xi_i=\frac{|Z_i|^2}{2}-\log n,\qquad \ell_n^{\mathrm{far}}=8\log\rho_n,
\]
and let
\[
    \mathcal R_n(A)=\{r:\rho_n/2\le r\le \rho_n-A\}\cup\{\rho_n+A\}.
\]
Completing the square gives the radial identity
\begin{equation}
\label{eq:d2-diffuse-radial-identity}
    \frac1n\exp\{rv\cdot z-r^2/2\}
    =
    \exp\left\{x(z)-\frac{|z-rv|^2}{2}\right\},
    \qquad
    x(z)=\frac{|z|^2}{2}-\log n .
\end{equation}

First remove the near-edge layer. Let \(H_n^{\mathrm{diff},0}\) be
\(H_n^{\mathrm{diff}}\) with the additional restriction
\(\xi_i\le -\ell_n^{\mathrm{far}}\). Every summand retained in
\(H_n^{\mathrm{diff},0}\) is at most
\(e^{-\ell_n^{\mathrm{far}}}\) after normalization. Hence the unnormalized summand is bounded
by \(ne^{-\ell_n^{\mathrm{far}}}\), and its second moment is at most \(ne^{-\ell_n^{\mathrm{far}}}\) times its
first moment.  Since
\[
    \E H_n^{\mathrm{diff},0}(r,v)
    \le \E H_n^{\mathrm{diff}}(r,v)=\Phi(\rho_n-r)\le1
    \qquad (r\in\mathcal R_n(A)),
\]
Bernstein's inequality gives, at each fixed \((r,v)\),
\[
    \Pbb\{|H_n^{\mathrm{diff},0}(r,v)
                 -\E H_n^{\mathrm{diff},0}(r,v)|>\eta\}
    \le 2\exp\{-c_\eta e^{\ell_n^{\mathrm{far}}}\}.
\]
Let \(\mathcal F_n\) be the class of normalized summands
\[
    z\mapsto \frac1n e^{rv\cdot z-r^2/2}
        \mathbf 1_{\{v\cdot z\le\rho_n\}}
        \mathbf 1_{\{x(z)\le -\ell_n^{\mathrm{far}}\}},
    \qquad r\in\mathcal R_n(A),\ v\in S^1,
\]
so that \(H_n^{\mathrm{diff},0}(r,v)\) is the sum of \(n\) functions from
\(\mathcal F_n\). This is a VC-subgraph class of fixed dimension: after taking
logarithms in the subgraph coordinate, its subgraphs are intersections of a
fixed ball, a halfspace in \(z\), and a halfspace in \((z,\log y)\).
Standard closure properties therefore give polynomial covering numbers,
uniformly in \(n\); see
\cite[Sections~2.6.2 and~2.14.1]{vanderVaartWellner1996}. In the present
bounded-envelope setting, the Bernstein maximal inequality from that entropy
bound has the form
\[
 \Pbb\left\{\sup_{f\in\mathcal F_n}
   \left|\sum_{i=1}^n\{f(Z_i)-\E f(Z_i)\}\right|>\eta\right\}
 \le C\rho_n^C
 \exp\left\{-\frac{c\eta^2}
 {e^{-\ell_n^{\mathrm{far}}}+\eta e^{-\ell_n^{\mathrm{far}}}}\right\}.
\]
Moreover, identity~\eqref{eq:d2-diffuse-radial-identity}
bounds every member of \(\mathcal F_n\) by \(e^{-\ell_n^{\mathrm{far}}}\), and
the sum of its variances by \(e^{-\ell_n^{\mathrm{far}}}\). The Bernstein maximal inequality for
VC-subgraph classes now upgrades the pointwise estimate to the whole class.
Its entropy cost is \(O(\log\rho_n)\), whereas its deviation exponent is of
order \(e^{\ell_n^{\mathrm{far}}}=\rho_n^8\). Hence
\[
    \sup_{r\in\mathcal R_n(A),\,v\in S^1}
    |H_n^{\mathrm{diff},0}(r,v)-\E H_n^{\mathrm{diff},0}(r,v)|
    \xrightarrow{\Pbb}0 .
\]

It remains to control the removed layer.  Write \(Z_i=R_iU_i\) and
\(Y_i(v)=\rho_n\operatorname{ang}(U_i,v)\).  If \(\xi_i>0\) and \(v\cdot Z_i\le\rho_n\),
then
\[
    |Z_i-rv|^2
    \ge |Z_i|^2+r^2-2r\rho_n
    =2\xi_i+(r-\rho_n)^2,
\]
so every such positive-edge term is bounded by \(e^{-A^2/2}\), uniformly
over \(r\in\mathcal R_n(A)\).  Thus their total contribution is at most
\(e^{-A^2/2}N_n^+\), where \(N_n^+=\#\{i:\xi_i>0\}=O_{\Pbb}(1)\).

For \(-\ell_n^{\mathrm{far}}<\xi_i\le0\), \(R_i=\rho_n+O(\ell_n^{\mathrm{far}}/\rho_n)\).  Uniformly over
\(r\in\mathcal R_n(A)\), for all large \(n\),
\[
    |Z_i-rv|^2
    \ge cA^2+cY_i(v)^2 .
\]
Therefore the negative part of the removed layer is bounded by
\[
    e^{-cA^2}
    \sup_{v\in S^1}
    \sum_{-\ell_n^{\mathrm{far}}<\xi_i\le0} e^{\xi_i}e^{-cY_i(v)^2}.
\]
The last supremum is tight.  Indeed, partition \((-\ell_n^{\mathrm{far}},0]\) into the slabs
\(-k-1<\xi_i\le-k\), \(0\le k\le\lceil \ell_n^{\mathrm{far}}\rceil\), and cover \(S^1\) by
\(O(\rho_n)\) arcs of length \(O(\rho_n^{-1})\).  If \(N_{k\ell}\) denotes
the number of points in the \(k\)-th slab and in the \(\ell\)-th enlarged
arc, then the Gaussian angular kernel has summable tails and
\[
    \sup_v\sum_{-\ell_n^{\mathrm{far}}<\xi_i\le0} e^{\xi_i}e^{-cY_i(v)^2}
    \le
    C\sum_{k\le\lceil \ell_n^{\mathrm{far}}\rceil} e^{-k}\max_\ell N_{k\ell}.
\]
Here \(\E N_{k\ell}\le C e^k/\rho_n\).  Chernoff's bound and a union bound
over the \(O(\rho_n)\) arcs give, with probability tending to one,
\[
    \max_\ell N_{k\ell}
    \le C\left(\frac{e^k}{\rho_n}+k+1\right)
    \qquad\text{for all }k\le\lceil \ell_n^{\mathrm{far}}\rceil .
\]
Consequently the displayed shot-noise supremum is \(O_{\Pbb}(1)\).  We have
proved
\[
    \sup_{r\in\mathcal R_n(A),\,v\in S^1}
    \{H_n^{\mathrm{diff}}(r,v)-H_n^{\mathrm{diff},0}(r,v)\}
    \le e^{-A^2/2}O_{\Pbb}(1)+e^{-cA^2}O_{\Pbb}(1).
\]

On the lower side, \(\E H_n^{\mathrm{diff},0}(r,v)\le1\), so the first
assertion follows
from the concentration estimate and then \(A\to\infty\).  At the upper
boundary,
\(\E H_n^{\mathrm{diff},0}(\rho_n+A,v)
\le\E H_n^{\mathrm{diff}}(\rho_n+A,v)=\Phi(-A)\), and
\(\Phi(-A)\to0\). The concentration and removed-layer bounds above therefore
prove the second assertion as well.
\end{proof}

The tail part \(T_n\) is likewise negligible away from
\(|r-\rho_n|=O(1)\).
\begin{lemma}
\label{lem:d-two-tail-outside}
For every \(\delta>0\),
\[
    \lim_{A\to\infty}\limsup_{n\to\infty}
    \Pbb\left\{
        \sup_{\substack{r\ge\rho_n/2,\ v\in S^1\\ |r-\rho_n|\ge A}}
        T_n(r,v)>\delta
    \right\}=0 .
\]
\end{lemma}

\begin{proof}
Put \(\xi_i=|Z_i|^2/2-\log n\) and \(N_n^+=\#\{i:\xi_i>0\}\).  Since
\(|Z_i|^2/2\) is exponential in dimension two, \(N_n^+\Rightarrow
\operatorname{Pois}(1)\) and
\[
    \Pbb\{\max_i\xi_i>A^2/4\}\to 1-\exp\{-e^{-A^2/4}\}.
\]
Hence the event
\[
    \max_i\xi_i\le A^2/4,\qquad N_n^+\le e^{A^2/8}
\]
has probability tending to one as first \(n\to\infty\) and then
\(A\to\infty\).  If \(v\cdot Z_i>\rho_n\), then \(\xi_i>0\).  On the
displayed event, write \(R_i=|Z_i|\).  For \(r\le\rho_n-A\),
\[
    r\,v\cdot Z_i-\frac{r^2}{2}-\log n
    \le
    (\rho_n-A)R_i-\frac{(\rho_n-A)^2}{2}-\log n
    \le
    \xi_i-\frac{A^2}{2}+o(1).
\]
For \(r\ge\rho_n+A\), the maximum over \(r\) is attained either at
\(r=\rho_n+A\) or at \(r=v\cdot Z_i\).  Since \(\xi_i\le A^2/4\) implies
\(R_i=\rho_n+O(A^2/\rho_n)<\rho_n+A\) for large \(n\), the boundary case
applies and gives
\[
    r\,v\cdot Z_i-\frac{r^2}{2}-\log n
    \le
    (\rho_n+A)R_i-\frac{(\rho_n+A)^2}{2}-\log n
    \le
    \xi_i-\frac{A^2}{2}+o(1)
    \le -A^2/4+o(1).
\]
Therefore
\[
    \sup_{\substack{r\ge\rho_n/2,\ v\in S^1\\ |r-\rho_n|\ge A}}
    T_n(r,v)
    \le
    N_n^+e^{-A^2/4+o(1)}
    \le e^{-A^2/8+o(1)}
\]
on the displayed event, which proves the lemma.
\end{proof}

An ordinary bulk estimate removes the remaining range
\(0<|t|\le\rho_n/2\).
\begin{lemma}
\label{lem:d2-bulk-exclusion}
Assume \(d=2\) and \(\lambda_n=\kappa_n\rho_n^2\to\lambda\in(0,\infty)\).
Then
\[
    \Pbb\left\{\sup_{0<|t|\le\rho_n/2}D_n(t)>1\right\}\to0 .
\]
\end{lemma}

\begin{proof}
We repeat the three estimates behind
Proposition~\ref{prop:bulk-mgf-exclusion}, proved for fixed \(d\ge3\) in
Section~\ref{sec:sharp-high-dimensional-estimates}, and verify them at the
two-dimensional scale
\(\kappa_n\asymp\rho_n^{-2}\asymp1/\log n\).

First consider \(0<|t|\le c_0\kappa_n\). The sample covariance satisfies
\[
    \left\|\frac1n\sum_i
    (Z_i-\bar Z_n)(Z_i-\bar Z_n)^\top-I_2\right\|_{\mathrm{op}}
    =O_{\Pbb}(n^{-1/2})
    =o_{\Pbb}(\kappa_n).
\]
The Hessian argument in Lemma~\ref{lem:microscopic-bulk} therefore applies
unchanged: for a sufficiently small fixed \(c_0>0\), with probability tending
to one,
\[
    D_n(t)\le1-c\kappa_n|t|^2<1
    \qquad(0<|t|\le c_0\kappa_n).
\]

For \(c_0\kappa_n\le |t|\le1\), use the quadratic empirical-process
expansion from Lemma~\ref{lem:bulk-exponential-process}. The only scale
conditions needed there are
\[
    \frac{n^{-1/4}}{\kappa_n}\to0,
    \qquad
    \frac{|\bar Z_n|}{\kappa_n^2}
    =
    O_{\Pbb}(n^{-1/2}\rho_n^4)\to0 .
\]
They imply, uniformly on this range,
\[
    \frac{|H_n(t)-1|+|t\cdot\bar Z_n|}
         {1-e^{-\kappa_n|t|^2}}
    =o_{\Pbb}(1).
\]

For \(1\le |t|\le\rho_n/2\), repeat the dyadic decomposition and clipping
argument from Lemma~\ref{lem:bulk-exponential-process}. Explicitly, partition the squared radii into
dyadic intervals \([x_j,x_{j+1}]\), put
\[
 R_j=x_{j+1}^{1/2},\qquad
 \tau_n^2=20\log(1/\kappa_n)+100\log\rho_n,qquad
 B_j=\exp\{R_j^2/2+R_j\tau_n\},
\]
and clip each likelihood ratio at \(B_j\). Here
\(|t|^2\le(\log n)/2\), so the variance and envelope bounds are
\(e^{|t|^2}\le n^{1/2}\) and \(B_j\le n^{1/4+o(1)}\). The net entropy is
still \(O(\log n)\), while the
Bernstein exponent is at least
\[
    n^{1/2-o(1)}\kappa_n^2
    \gg \log n .
\]
Thus there is a deterministic \(\delta_n\downarrow0\) such that, with
probability tending to one,
\[
    |H_n(t)-1|+|t\cdot\bar Z_n|
    \le\delta_n g(t),
    \qquad
    g(t)=1-e^{-\kappa_n|t|^2},
\]
uniformly on \(c_0\kappa_n\le|t|\le\rho_n/2\). (In particular,
\(\rho_n|\bar Z_n|=o_{\Pbb}(1)\).) The algebra in the proof of
Proposition~\ref{prop:bulk-mgf-exclusion} then gives
\[
    D_n(t)
    \le e^{\delta_ng(t)}\{1-g(t)\}\{1+\delta_ng(t)\}
    <1
\]
uniformly on that range for all large \(n\). Together with the microscopic
estimate, this proves the lemma.
\end{proof}

\begin{proof}[Proof of Proposition~\ref{prop:d-two-radial-compactness}]
Lemma~\ref{lem:d2-bulk-exclusion} controls the bulk range:
\[
    \Pbb\left\{\sup_{0<|t|\le\rho_n/2}D_n(t)>1\right\}\to0 .
\]

For the upper exterior, use the centered representation
\[
    D_n(rv)=\frac1n\sum_i
    \exp\left\{rv\cdot (Z_i-\bar Z_n)-\frac{1+2\kappa_n}{2}r^2\right\}.
\]
With probability tending to one, \(\max_i|Z_i-\bar Z_n|<\rho_n+A/2\).
On this event every summand is decreasing in \(r\) for \(r\ge\rho_n+A\),
so the upper exterior reduces to the boundary \(r=\rho_n+A\).  Since
\(\rho_n|\bar Z_n|=o_{\Pbb}(1)\),
\[
    \sup_{v}D_n((\rho_n+A)v)
    \le
    e^{-\lambda+o_{\Pbb}(1)}
    \sup_v\{H_n^{\mathrm{diff}}(\rho_n+A,v)+T_n(\rho_n+A,v)\}.
\]
The upper-bound parts of Lemmas~\ref{lem:d-two-diffuse-outside} and
\ref{lem:d-two-tail-outside} make the right side smaller than one with
probability tending to one after \(A\to\infty\).

It remains to consider \(\rho_n/2\le r\le\rho_n-A\).  On this range
\[
    e^{-\kappa_nr^2}\le e^{-\lambda/4+o(1)},
    \qquad
    \sup_{r,v}e^{-rv\cdot\bar Z_n}=e^{o_{\Pbb}(1)} .
\]
Lemma~\ref{lem:d-two-diffuse-outside} gives
\(H_n^{\mathrm{diff}}(r,v)\le1+o_{\Pbb}(1)+o_A(1)\) uniformly, and
Lemma~\ref{lem:d-two-tail-outside} gives
\(T_n(r,v)=o_{\Pbb}(1)+o_A(1)\) uniformly.  Since \(e^{-\lambda/4}<1\),
the lower exterior also cannot contain a crossing once \(A\) is large.
\end{proof}

\subsection{Compact-window decomposition}

We now approximate the full field on the compact radial window. Observations
below a slowly receding radial cutoff form the
deterministic background; those above it produce narrow angular kernels.
\begin{proposition}
\label{prop:two-d-compact-decomposition}
Assume \(d=2\), put \(\rho_n=\sqrt{2\log n}\), and let \(K\subset\R\) be compact.  Let \(\operatorname{ang}(u,v)\in(-\pi,\pi]\) denote the principal signed angular displacement from \(v\) to \(u\).  For \(Z_i\ne0\), write \(Z_i=R_iU_i\), and set
\[
    \xi_i=\frac{R_i^2}{2}-\log n,
    \qquad
    Y_i(v)=\rho_n\,\operatorname{ang}(U_i,v),
\]
so that \(Y_i(v)\) is the blown-up signed angular displacement from \(v\) to
the observation direction \(U_i\). If \(I\subset S^1\) is a closed arc and
\(\ell_n^{\mathrm{edge}}=4\log\rho_n\), then
\[
    \sup_{v\in I,\ s\in K}
    \left|
    H_n((\rho_n+s)v)-\Phi(-s)
    -
    \sum_{i:\xi_i\ge -\ell_n^{\mathrm{edge}}}
    \exp\left\{\xi_i-\frac{s^2}{2}-\frac{Y_i(v)^2}{2}\right\}
    \right|
    \xrightarrow{\Pbb}0 .
\]
\end{proposition}

\begin{proof}
Write \(r_s=\rho_n+s\).  First separate the diffuse contribution below the
radial edge.  For \(L>0\), define
\[
    H^{\rm diff}_{n,L}(v,s)
    =
    \frac1n\sum_{i=1}^n
    \exp\left\{r_sv\cdot Z_i-\frac{r_s^2}{2}\right\}
    \mathbf 1_{\{\xi_i<-L\}} .
\]
We use \(L=\ell_n^{\mathrm{edge}}\).  If \(\xi_i<-L\), then \(R_i\le R_L=(\rho_n^2-2L)^{1/2}\), and uniformly for \(s\in K\),
\[
    \frac1n\exp\left\{r_sv\cdot Z_i-\frac{r_s^2}{2}\right\}
    \le
    \exp\left\{-L-\frac{(r_s-R_L)^2}{2}\right\}
    \le C_K e^{-L}.
\]
The sum of the variances is therefore at most \(C_Ke^{-L}\).  Bernstein's inequality gives, for each fixed \((v,s)\),
\[
    \Pbb\{|H^{\rm diff}_{n,L}(v,s)-\E H^{\rm diff}_{n,L}(v,s)|>\eta\}
    \le
    2\exp\{-c_{K,\eta}e^L\}.
\]
The field \(H^{\rm diff}_{n,L}\) is smooth in \((v,s)\). Each summand has
logarithmic derivatives bounded by \(C_K\rho_n^2\) in the angular coordinate
and by \(C_K\rho_n\) in \(s\). Fix \(\eta>0\), and choose an angular mesh of
order \(\rho_n^{-2}\) and an \(s\)-mesh of order \(\rho_n^{-1}\). The mesh
constants can be chosen so that neighboring values of both
\(H^{\rm diff}_{n,L}\) and \(\E H^{\rm diff}_{n,L}\) differ by a factor at
most \(1+\eta\). The grid has polynomial size in \(\rho_n\), so a union bound
upgrades the pointwise Bernstein estimate to the grid. Because
\(\E H^{\rm diff}_{n,L}\le1\), interpolation from the grid costs only an
additive \(O(\eta)\); letting \(\eta\downarrow0\) gives
\[
    \sup_{v\in I,\ s\in K}
    |H^{\rm diff}_{n,\ell_n^{\mathrm{edge}}}(v,s)-\E H^{\rm diff}_{n,\ell_n^{\mathrm{edge}}}(v,s)|
    \xrightarrow{\Pbb}0,
\]
because \(e^{\ell_n^{\mathrm{edge}}}=\rho_n^4\) dominates the grid entropy.

It remains to compute the mean.  By exponential tilting,
\[
    \E H^{\rm diff}_{n,L}(v,s)
    =
    \Pbb\{|G+r_sv|^2\le \rho_n^2-2L\},
    \qquad G\sim\cN(0,I_2).
\]
After rotating \(v=e_1\), this probability is
\[
    \Pbb\left\{(r_s+G_1)^2+G_2^2\le \rho_n^2-2L\right\}.
\]
Condition on \(G_2\).  On the event \(|G_2|\le\rho_n^{1/3}\), the boundary for \(G_1\) is
\[
    G_1
    \le
    -s-\frac{\ell_n^{\mathrm{edge}}+G_2^2/2}{\rho_n}
    +O_K\left(\frac{(\ell_n^{\mathrm{edge}}+G_2^2+1)^2}{\rho_n^3}\right),
\]
uniformly for \(s\in K\). The quadratic inequality also has the lower
boundary
\[
 G_1\ge-r_s-\sqrt{\rho_n^2-2\ell_n^{\mathrm{edge}}-G_2^2}
       =-2\rho_n+O_K(1)
\]
on this event, whose Gaussian probability is uniformly \(o(1)\). The
complement of \(\{|G_2|\le\rho_n^{1/3}\}\) also has probability \(o(1)\),
and the displayed upper boundary converges to \(-s\) uniformly on \(K\)
after conditioning on \(G_2\). Dominated convergence and the Lipschitz
continuity of the Gaussian distribution function give
\[
    \sup_{v\in I,\ s\in K}
    |\E H^{\rm diff}_{n,\ell_n^{\mathrm{edge}}}(v,s)-\Phi(-s)|\to0 .
\]
Thus the contribution from \(\{\xi_i<-\ell_n^{\mathrm{edge}}\}\) is \(\Phi(-s)+o_{\Pbb}(1)\), uniformly on \(I\times K\).

Now consider the edge contribution
\[
    A^{\rm ex}_n(v,s)
    =
    \sum_{i:\xi_i\ge -\ell_n^{\mathrm{edge}}}
    \frac1n
    \exp\left\{r_sv\cdot Z_i-\frac{r_s^2}{2}\right\}.
\]
We compare it with the displayed atom sum. Let
\[
    C_n=\log\rho_n,
    \qquad
    Y_n=(\log\rho_n)^{1/4}.
\]
The event \(\max_i\xi_i\le C_n\) has probability tending to one, since
\(n\Pbb\{\xi_i>C_n\}=e^{-C_n}\). Also
\[
    \sum_{i:-\ell_n^{\mathrm{edge}}\le \xi_i\le C_n} e^{\xi_i}
    =
    O_{\Pbb}(\ell_n^{\mathrm{edge}}+C_n),
\]
because \(\xi_i\) has shifted exponential intensity \(e^{-x}\,dx\) in dimension two.

On the event \(\max_i\xi_i\le C_n\), the exact contribution of an edge observation
with \(-\ell_n^{\mathrm{edge}}\le \xi_i\le C_n\) satisfies
\[
    \frac1n
    \exp\left\{r_sv\cdot Z_i-\frac{r_s^2}{2}\right\}
    \le
    C_K e^{\xi_i}
    \exp\{-c\min(Y_i(v)^2,\rho_n^2)\}.
\]
This follows from
\[
    r_sv\cdot Z_i-\frac{r_s^2}{2}-\log n
    =
    \xi_i-\frac{(R_i-r_s)^2}{2}-r_sR_i\{1-\cos\operatorname{ang}(U_i,v)\},
\]
and \(r_sR_i\asymp\rho_n^2\) in the displayed range.  Hence the total
contribution of terms with \(|Y_i(v)|>Y_n\), in either \(A^{\rm ex}_n\) or
the Gaussian atom sum, is
\[
    O_{\Pbb}(\ell_n^{\mathrm{edge}}+C_n)e^{-cY_n^2}=o_{\Pbb}(1)
\]
uniformly in \(v\in I\) and \(s\in K\), by the definition of \(Y_n\).
For the remaining terms
\(-\ell_n^{\mathrm{edge}}\le \xi_i\le C_n\) and \(|Y_i(v)|\le Y_n\), Taylor expansion gives,
uniformly in \(i,v,s\),
\[
\begin{aligned}
    &\log\left[
    \frac1n\exp\left\{r_sv\cdot Z_i-\frac{r_s^2}{2}\right\}
    \right]
    -\left\{\xi_i-\frac{s^2}{2}-\frac{Y_i(v)^2}{2}\right\} \\
    &\qquad=
    O_K\left(\frac{\ell_n^{\mathrm{edge}}+C_n}{\rho_n}
    +\frac{Y_n^2}{\rho_n}
    +\frac{Y_n^4}{\rho_n^2}\right)
    =O_K\left(\frac{\log\rho_n}{\rho_n}\right).
\end{aligned}
\]
The total atom weight is \(O_{\Pbb}(\log\rho_n)\), while the relative error is
\(O(\log\rho_n/\rho_n)\). Their product is \(o_{\Pbb}(1)\), and hence
\[
    \sup_{v\in I,\ s\in K}
    \left|
    A^{\rm ex}_n(v,s)
    -
    \sum_{i:\xi_i\ge -\ell_n^{\mathrm{edge}}}
    \exp\left\{\xi_i-\frac{s^2}{2}-\frac{Y_i(v)^2}{2}\right\}
    \right|
    \xrightarrow{\Pbb}0 .
\]
Combining this with \(H_n=H^{\rm diff}_{n,\ell_n^{\mathrm{edge}}}+A^{\rm ex}_n\) proves the proposition.
\end{proof}

Proposition~\ref{prop:two-d-compact-decomposition} uses the slowly receding
cutoff \(\ell_n^{\mathrm{edge}}=4\log\rho_n\), which makes the background concentrate uniformly.
For the final supremum bound, we replace it by a fixed cutoff. For every fixed
\(A<\infty\), the
edge points with \(\xi_i>-A\) are finite in the limit and have no
\(1/\rho_n\)-scale
angular collisions:
\[
    \Pbb\left\{
    \exists i\ne j:\ \xi_i,\xi_j>-A,\
    \rho_n|\operatorname{ang}(U_i,U_j)|\le L
    \right\}\le C_A L/\rho_n\to0 .
\]

Edge points in the lower radial strip \(-\ell_n^{\mathrm{edge}}\le \xi_i\le-A\) have negligible
total contribution after \(A\to\infty\), allowing a fixed cutoff.
\begin{lemma}
\label{lem:d2-low-atom-absorption}
Let \(\ell_n^{\mathrm{edge}}=4\log\rho_n\), and write
\(Y_i(v)=\rho_n\operatorname{ang}(U_i,v)\) for the principal angular
displacement between \(U_i\) and \(v\).  For every
\(\delta>0\),
\[
    \lim_{A\to\infty}\limsup_{n\to\infty}
    \Pbb\left\{
    \sup_{v\in S^1}
    \sum_{i:-\ell_n^{\mathrm{edge}}\le \xi_i\le -A}
    e^{\xi_i-Y_i(v)^2/2}>\delta
    \right\}=0 .
\]
\end{lemma}

\begin{proof}
Put \(m_n^\star=\lceil \ell_n^{\mathrm{edge}}\rceil\) and, for integers \(m\ge0\),
\[
    I_{m,n}=\{i:-m-1<\xi_i\le -m\}.
\]
These are unit radial-excess layers; on \(I_{m,n}\), the weight \(e^{\xi_i}\)
is of order \(e^{-m}\).  Since
\(|Z_i|^2/2\) is exponential in dimension two,
\[
    \Pbb\{i\in I_{m,n}\}\le C e^m/n,
    \qquad m\le m_n^\star,
\]
and the angle \(\Theta_i\) is independent and uniform on the circle.

For \(q=0,1,\ldots,\lfloor\pi\rho_n\rfloor\), let
\[
    N_{m,q,n}
    =
    \sup_{\phi\in\mathbb T}
    \#\left\{
        i\in I_{m,n}:
        \rho_n|\operatorname{ang}(\Theta_i-\phi)|\le q+1
    \right\}.
\]
Covering the circle by \(O(\rho_n)\) arcs shows that every interval appearing
in the supremum is contained in one of the covering arcs, with length enlarged
by only a fixed factor.  For one such enlarged arc the count is binomial with
mean
\[
    \mu_{m,q,n}\le C(q+1)e^m/\rho_n .
\]
Set
\[
    a_{m,q,n}
    =
    1+\frac{(q+1)e^m}{\rho_n}
    +
    \frac{\log\rho_n}
        {1+\log_+\!\left(\rho_n/((q+1)e^m)\right)},
    \qquad
    \log_+ x=\max\{\log x,0\}.
\]
The binomial Chernoff bound
\(\Pbb\{\operatorname{Bin}(N,p)\ge r\}\le (eNp/r)^r\) for
\(r\ge eNp\), together with the multiplicative Chernoff bound when
\(r\) is a fixed multiple of \(Np\), implies that for a sufficiently large
numerical constant \(B\),
\[
    \Pbb\left\{
        N_{m,q,n}>B a_{m,q,n}
        \text{ for some } A\le m\le m_n^\star,\ 0\le q\le \pi\rho_n
    \right\}\to0
\]
for each fixed \(A\).  Indeed, if \(\mu_{m,q,n}<1\), the last term in
\(a_{m,q,n}\) makes
\[
    a_{m,q,n}\log\{a_{m,q,n}/\mu_{m,q,n}\}\ge c\log\rho_n,
\]
while if \(\mu_{m,q,n}\ge1\), the term \(\log\rho_n\) gives the same
union-bound margin unless \(\mu_{m,q,n}\ge\log\rho_n\), in which case the
linear mean term does.  Choosing \(B\) large makes the tail summable over
the \(O(\rho_n^2\log\rho_n)\) enlarged arcs, layers, and angular scales.

On this high-probability occupancy event, uniformly in \(v\),
\[
\begin{aligned}
    \sum_{i:-\ell_n^{\mathrm{edge}}\le \xi_i\le-A}e^{\xi_i-Y_i(v)^2/2}
    &\le
    C\sum_{m=\lfloor A\rfloor}^{m_n^\star} e^{-m}
    \sum_{q=0}^{\lfloor\pi\rho_n\rfloor}
        e^{-q^2/2}N_{m,q,n} \\
    &\le
    C\sum_{q=0}^{\lfloor\pi\rho_n\rfloor}e^{-q^2/2}
    \sum_{m=\lfloor A\rfloor}^{m_n^\star} e^{-m}a_{m,q,n}.
\end{aligned}
\]
The first two terms in \(a_{m,q,n}\) contribute
\[
    O(e^{-A})+O(\ell_n^{\mathrm{edge}}/\rho_n).
\]
For the logarithmic occupancy term, fix \(A\) and split the \(m\)-sum at
\(\log\{\rho_n/(q+1)\}\).  This gives the deterministic bound
\[
    \sum_{m=\lfloor A\rfloor}^{m_n^\star}
    e^{-m}
    \frac{\log\rho_n}
        {1+\log_+\!\left(\rho_n/((q+1)e^m)\right)}
    \le
    C e^{-A}\{1+\log(q+2)\}
    +C\frac{(q+1)\log\rho_n}{\rho_n}
\]
for all \(n\ge n_0(A)\), uniformly in \(q\). This fixed-\(A\) qualification
is essential; the estimate is not asserted for a simultaneous choice
\(A=A_n\).
After summing in \(q\) against \(e^{-q^2/2}\), the right side is
\[
    O(e^{-A})+o(1).
\]
Thus the displayed supremum is at most \(C e^{-A}+o(1)\) with probability
tending to one, and the lemma follows by letting \(A\to\infty\).
\end{proof}

The moving-cutoff decomposition can therefore be replaced by the fixed-cutoff
form needed for the reverse inequality.
\begin{lemma}
\label{lem:d2-fixed-cutoff-decomposition}
For each compact \(K\subset\R\), each \(C<\infty\), and each \(\delta>0\),
\[
\begin{aligned}
&\lim_{A\to\infty}\limsup_{n\to\infty}
\Pbb\Biggl\{
\max_i\xi_i\le C,\
\sup_{\substack{v\in S^1\\ s\in K}}
\Biggl|
H_n((\rho_n+s)v)-\Phi(-s)\\
&\hspace{35mm}
-\sum_{i:-A<\xi_i\le C}
e^{\xi_i-s^2/2-Y_i(v)^2/2}
\Biggr|>\delta
\Biggr\}=0 .
\end{aligned}
\]
\end{lemma}

\begin{proof}
Apply Proposition~\ref{prop:two-d-compact-decomposition} on a finite closed-arc
cover of \(S^1\). On \(\{\max_i\xi_i\le C\}\), its fixed-cutoff sum omits only
the observations with
\(-\ell_n^{\mathrm{edge}}\le \xi_i\le-A\). Their
absolute contribution is bounded uniformly in \(s\in K\) by
\[
    \sup_{v\in S^1}
    \sum_{i:-\ell_n^{\mathrm{edge}}\le \xi_i\le-A} e^{\xi_i-Y_i(v)^2/2},
\]
because \(e^{-s^2/2}\le1\).  Lemma~\ref{lem:d2-low-atom-absorption} makes
this term vanish after first taking \(n\to\infty\) and then \(A\to\infty\),
proving the fixed-cutoff form.
\end{proof}
